\documentclass[tikz,border=8pt]{standalone}
\usepackage{tikz}
\usepackage{amsmath}
\usepackage{amssymb}
\usepackage{pgfplots}
\pgfplotsset{compat=1.18}
\usepackage{ctex}
\usetikzlibrary{calc,positioning,arrows.meta,matrix,trees}

% Trie 字典树结构
% 插入 "cat", "car", "card" 后的 Trie
% 根 → c → a → t(end) / r → end / r → d(end)

\begin{document}
\begin{tikzpicture}[
  every node/.style={font=\small},
  trienode/.style={draw=blue!60, circle, minimum size=0.75cm, thick, fill=blue!8},
  endnode/.style={draw=green!60!black, circle, minimum size=0.75cm, very thick, fill=green!20},
  rootnode/.style={draw=gray!70, circle, minimum size=0.75cm, thick, fill=gray!10},
  edgelabel/.style={font=\small\bfseries, midway, fill=white, inner sep=1pt},
  ptr/.style={->,>=Stealth,thick,draw=gray!70},
]

%% 标题
\node[font=\large\bfseries] at (4.5, 1.5)
  {Trie 字典树（插入 "cat", "car", "card"）};
\node[font=\small,gray] at (4.5, 0.9)
  {绿色节点 = isEnd，蓝色 = 中间节点，边上字母为字符};

%% Trie 结构
%% Level 0: root
%% Level 1: c
%% Level 2: a
%% Level 3: t (isEnd=cat), r (isEnd=car)
%% Level 4: d (isEnd=card，挂在 r 下)

% root
\node[rootnode] (root) at (4.5, 0.2) {root};

% c
\node[trienode] (c) at (4.5, -1.2) {c};
\draw[ptr] (root) -- (c) node[edgelabel] {c};

% a
\node[trienode] (a) at (4.5, -2.7) {a};
\draw[ptr] (c) -- (a) node[edgelabel] {a};

% t (cat end) and r (car end)
\node[endnode] (t) at (2.2, -4.2) {t};
\node[endnode] (r) at (6.8, -4.2) {r};
\draw[ptr] (a) -- (t) node[edgelabel] {t};
\draw[ptr] (a) -- (r) node[edgelabel] {r};

% d (card end, hangs under r)
\node[endnode] (d) at (6.8, -5.7) {d};
\draw[ptr] (r) -- (d) node[edgelabel] {d};

%% 单词终结标注（放在节点旁，不遮挡）
\node[font=\scriptsize,green!60!black,anchor=west] at (2.65, -4.2) {← "cat" 结束};
\node[font=\scriptsize,green!60!black,anchor=west] at (7.25, -4.2) {← "car" 结束};
\node[font=\scriptsize,green!60!black,anchor=west] at (7.25, -5.7) {← "card" 结束};

%% isEnd 图例
\node[font=\small\bfseries] at (1.0, -1.5) {图例};
\node[endnode,minimum size=0.55cm]  at (0.6, -2.2) {};
\node[font=\scriptsize,anchor=west] at (0.95, -2.2) {isEnd=true};
\node[trienode,minimum size=0.55cm] at (0.6, -2.9) {};
\node[font=\scriptsize,anchor=west] at (0.95, -2.9) {isEnd=false};

%% Trie 节点内部结构说明（右侧）
\node[draw=gray!40,fill=white,thin,rounded corners=2pt,
      font=\scriptsize,inner sep=6pt,align=left]
  at (1.5, -6.8)
  {\textbf{TrieNode 结构：}\\[3pt]
   children: dict \{ char $\to$ TrieNode \}\\[2pt]
   isEnd: bool（该节点为某词结尾）};

%% 插入过程说明
\node[draw=gray!40,fill=white,thin,rounded corners=2pt,
      font=\scriptsize,inner sep=6pt,align=left]
  at (6.5, -6.8)
  {\textbf{插入 "card" 流程：}\\[3pt]
   root $\to$ c（已有）$\to$ a（已有）\\[2pt]
   $\to$ r（已有，isEnd=true）\\[2pt]
   $\to$ d（新建节点，isEnd=true）};

%% 算法说明
\node[draw=gray!50,fill=white,rounded corners=3pt,font=\small,
      inner sep=8pt,align=left]
  at (4.5, -8.8)
  {\textbf{Trie 字典树：O(L) 插入/查询（L=单词长度）}\\[4pt]
   插入：从根出发，按字符逐层查找/新建节点，最后标 isEnd\\[2pt]
   查询：同插入路径，检查是否到达 isEnd 节点\\[2pt]
   前缀查询：只需遍历到前缀末尾，节点存在即可};

\end{tikzpicture}
\end{document}
